\documentclass{cicc}
\firstpage{1}
\articletype{Regular / Feature / Perspective / Review Article}

\doi{doi: 10.4208/cicc.2026.xxx.xx}
\publishedyear{2026}
\volume{1}
\issue{1}
%\pagenumbers{1-2}

\receiveddate{20 Oct. 2026}
\revisiondate{25 Oct. 2026}
\accepteddate{25 Dec. 2026}
\onlinedate{28 Dec. 2026}
\publisheddate{30 Dec. 2026}

\title[Here is the Title]
{Here is the Title}
\author[1,2]{John Doe}
\author[1,2,*]{Jane Smith}

\affil[1]{\textit{Department of Chemistry, Faculty of Arts and Sciences, Beijing Normal University, Zhuhai 519087, China} }
\affil[2]{\textit{College of Chemistry, Beijing Normal University, Beijing 100875, China} 
\protect\vspace{1em}}

\affil[*]{Corresponding authors: xxxxx@xxx}


\let\leqslant=\leq

\newtheorem{theorem}{Theorem}[section]

\begin{document}

\twocolumn[{
  \vspace*{0.5em}
  \maketitle
  \thispagestyle{firstpage}
  \label{firstpage}

\begin{abstract}
The abstract should be a single paragraph that summarises the content of the article.  

\keywords{keyword1, keyword2, keyword3, keyword4.}
\end{abstract}

}] % End of twocolumn header block


\section{Introduction}
The introduction section should appear here\cite{berger1989local,tan_simple_moving,mckechnie2017arxiv,boor1973good,tullius1996comprehensive,authors_year_dataset,sheldrick1996shelxs,tullius1991title}.

\section{Theoretical and computational details}
The theoretical and computational methods or details should appear here.

\subsection{Methods}
The time-dependent Schr\"odinger equation is:
\begin{equation} \label{eq:tdse}
i\hbar\frac{\partial\Psi(\mathbf{r},t)}{\partial t} = \hat{H}\Psi(\mathbf{r},t).
\end{equation}

\section{Results and discussion}
The results and discussions should appear here.

\subsection{Title of subsection}
All figures and tables must be mentioned in the text consecutively and numbered with Arabic numerals, such as Figure 1 or Table 1.

\begin{figure}[hbt!]
   \centering
   \includegraphics[width=1.0\linewidth]{Figure_1.jpg}
   \caption{Each figure must have a caption that includes the figure number and a brief description, preferably one or two sentences.}\label{sample-figure}
\end{figure}

% change \begin{table} to \begin{widetable} if you want to use the wide table environment
\begin{table}[hbt!]
\centering
\caption{Each table must have a brief (one phrase or sentence) title that describes its contents.}\label{sample-table}
\begin{tabular}{cccc}
\hline
Head 1\textsuperscript{[a]} & Head 2 & Head 3 & Head 4 \\
\hline
Column 1 & Column 2\textsuperscript{[b]} & Column 3 & Column 4 \\
Column 1 & Column 2 & Column 3 & Column 4 \\
Column 1 & Column 2 & Column 3 & Column 4 \\
Column 1 & Column 2 & Column 3 & Column 4 \\
Column 1 & Column 2 & Column 3 & Column 4 \\
\hline
\end{tabular}

\vspace{0.5em}
\raggedright\footnotesize
{[a] Table footnote. [b] Table footnote.}
\end{table}

\section{Conclusion}
The conclusion section should appear here.

\begin{acknowledgments}
Here is the place to acknowledge people, organizations, and financing (you may state grant numbers and sponsors here).
\end{acknowledgments}

\begin{Supporting_Information}
The online version contains supplementary material available at website xxx. A listing of the contents of each file supplied as Supporting Information should be included.
\end{Supporting_Information}

\begin{Notes}
The authors declare no competing financial interest.
\end{Notes}

\bibliographystyle{cicc}
\bibliography{ref}
\label{lastpage}
\end{document}

\section{Additional facilities}

\subsection{Titles and author's name}

In the CiCC style, the title of the article and the author's name (or
authors' names) are used both at the beginning of the article for the
main title and throughout the article as running headlines at the top
of every page. The title is used on odd-numbered pages (rectos) and the
author's name appears on even-numbered pages (versos). Although the
main heading can run to several lines of text, the running headline
must be a single line ($\leqslant 45$ characters). Moreover, the main
heading can also incorporate new line commands (e.g. \verb"\\") but
these are not acceptable in a running headline. To enable you to
specify an alternative short title and an alternative short author's
name, the standard \verb"\title" and \verb"\author" commands have been
extended to take an optional argument to be used as the running
headline. The running headlines for this guide were produced using the
following code:
\begin{verbatim}
\title[Commun. Comput. Chem.,
       \LaTeXe\ Guide for Authors]
\end{verbatim}
and
\begin{verbatim}
\author[B.L.N. Kennett]
   {B.L.N. Kennett$^1$
  \thanks{Pacific Region Office, CiCC} \\
  $^{1}$Research School of Earth Sciences,
    Australian National University,
    Canberra ACT \emph{0200}, Australia
  }
\end{verbatim}
The \verb"\thanks" note produces a footnote to the title or author.


\subsection{Captions for continued figures and tables}

The \verb"\caption*" command may be used to produce a caption with the
same number as the previous caption (for the corresponding type of
float). For instance, if a very large table does not fit on one page,
it must be split into two floats; the second float should use the
\verb"caption*" command with a suitable caption:
\begin{verbatim}
\begin{table}
 \caption*{-- \textit{continued}}
  \begin{tabular}{@{}lccll}
  :
  \end{tabular}
\end{table}
\end{verbatim}

 \begin{figure}
  \centering
  \includegraphics[width=0.3\textwidth]{CICC.jpg}
     \caption{An example figure in which space has been
              left for the artwork.}
     \label{sample-figure2}
  \end{figure}


\subsection{Tables}

The CiCC style will cope with positioning of your tables and you
should not use the positional qualifiers on the
\verb"table" environment which would override these decisions. Table
captions should be at the top, therefore the \verb"\caption" command
should appear before the body of the table.

The \verb"tabular" environment can be used to produce tables with
single horizontal rules, which are allowed, if desired, at the head and
foot only. This environment has been modified for the CiCC style in the
following ways:
\begin{enumerate}
  \item additional vertical space is inserted on either side of a rule;
  \item vertical lines are not produced.
\end{enumerate}
Commands to redefine quantities such as \verb"\arraystretch" should be
omitted. For example, Table~\ref{symbols} is produced using the
following commands.


\begin{table}
 \caption{Seismic velocities at major discontinuities.}
 \label{symbols}
 \begin{center}
 \begin{tabular}{||lcccccc||}
  \hline
  Class & depth & radius
        & $\alpha _{-}$ & $\alpha _{+}$
        & $\beta _{-}$ & $\beta _{+}$ \\
  \hline
  ICB & 5154 & 1217 & 11.091 & 10.258
        & 3.438 &  0. \\
  CMB & 2889 & 3482 & 8.009 & 13.691
        & 0. & 7.301 \\
  \hline
 \end{tabular}
 \end{center}

 \medskip

\end{table}


If you have a table that is to extend over two columns, you need to use
\verb"table*" in a minipage environment, i.e., you can say
\begin{verbatim}
\begin{table*}
\begin{minipage}{80mm}
 \caption{Caption which will wrap round to the
          width of the minipage environment.}
 \begin{tabular}{%
      :
 \end{tabular}
\end{minipage}
\end{table*}
\end{verbatim}
The width of the minipage should more or less be the width of your table,
so you can only guess on a value on the first pass. The value will have to
be adjusted when your article is finally typeset, so don't worry
about making it the exact size.

\subsection{Running headlines}

As described above, the title of the article and the author's name (or
authors' names) are used as running headlines at the top of every page.
The headline on right pages can list up to three names; for more than
three use et~al. The \verb"\pagestyle" and \verb"\thispagestyle"
commands should {\em not\/} be used. Similarly, t··he commands
\verb"\markright" and \verb"\markboth" should not be necessary.

\subsection{Bibliography}

The bibliography is generated automatically by BibTeX using the \verb"cicc.bst" style. Below is a summary of the formatting conventions for each entry type.

\subsubsection{General Rules}

\textbf{Author Names:} Surname first, followed by initials with no comma between surname and initials. Multiple initials are separated by periods and tied with non-breaking spaces. Multiple authors are separated by commas; the last author is preceded by ``and'' (no Oxford comma for two authors; ``and'' for three or more).\\
Example: \verb"Berger M.~J. and Colella P." or \verb"Tan Z.~J., Tang T. and Zhang Z.~R."

\textbf{Punctuation Between Fields:} Fields are separated by commas. Each entry ends with a period.

\subsubsection{Journal Articles (\texttt{@article})}

Format: Authors, Title, \emph{Journal Name}, \textbf{Volume}(Issue) (Year), Pages.

\begin{itemize}
  \item The article title appears in plain text (no italics, no quotes).
  \item The journal name is italicized and abbreviated. If a \verb"shortjournal" field is provided, it is used instead of the full journal name. Common full journal names in the output \verb".bib" copy should be normalized to abbreviations before BibTeX.
  \item The volume number is in bold. If an issue number is present, it follows in parentheses immediately after the volume.
  \item The year is enclosed in parentheses after the volume/issue.
  \item Page ranges follow the year, separated by a comma, using an en-dash.
\end{itemize}
Example output: Berger M.~J. and Colella P., Local adaptive mesh refinement for shock hydrodynamics, \emph{J. Comput. Phys.}, \textbf{82} (1989), 62--84.

\subsubsection{Books and Book Chapters (\texttt{@incollection}, \texttt{@inbook})}

Format: Authors, Title, In \emph{Book Title}, Vol.~\textbf{Volume}, (Eds.\ Editor Names), Publisher, Address, Year, pp.~Pages.

\begin{itemize}
  \item The chapter/article title appears in plain text.
  \item The book title is preceded by ``In'' and italicized.
  \item If a volume is present, it appears as ``Vol.~\textbf{N}''.
  \item Editors are listed in parentheses as ``(Eds.\ ...)'' or ``(Ed.\ ...)''.
  \item Publisher, address, and year follow, separated by commas.
  \item Page ranges are prefixed with ``pp.'' (or ``p.'' for a single page).
\end{itemize}
Example output: Tullius T.~D., Supramolecular reactivity and transport: Bioinorganic systems, In \emph{Comprehensive Supramolecular Chemistry}, Vol.~\textbf{5}, (Eds. Atwood J.~L., Davies J.~E.~D., MacNicol D.~D., V{\"o}gtle F. and Suslick K.~S.), Pergamon, Oxford, 1996, pp.~317--343.

\subsubsection{Standalone Books (\texttt{@book})}

Format: Authors, \emph{Book Title}, Publisher, Address, Year.

\begin{itemize}
  \item The book title is italicized.
  \item Publisher, address, and year follow, separated by commas.
\end{itemize}

\subsubsection{Preprints and Miscellaneous (\texttt{@misc})}

Format: Authors, Title, Year.

\begin{itemize}
  \item The title appears in plain text.
  \item DOI, URL, and ISSN fields are not displayed in the final reference list.
\end{itemize}
Example output: McKechnie S., Frost J.~M., ..., Dynamic symmetry breaking and spin splitting in metal halide perovskites, 2017.

\subsubsection{Datasets (\texttt{@dataset})}

Format: Authors, \emph{Dataset Title}, Year.

\begin{itemize}
  \item The dataset title is italicized (treated like a book title).
\end{itemize}

\subsubsection{Programs and Manuals (\texttt{@manual})}

Format: Authors, \emph{Program Title}, Organization, Address, Year.

\begin{itemize}
  \item The program/manual title is italicized.
  \item Organization and address follow, separated by commas.
\end{itemize}
Example output: Sheldrick G.~M., \emph{SHELXS-96, program for the solution of crystal structures}, University of G{\"o}ttingen, G{\"o}ttingen, Germany, 1996.

\subsubsection{Theses (\texttt{@phdthesis}, \texttt{@mastersthesis})}

Format: Authors, \emph{Thesis Title}, School, Address, Year.

\begin{itemize}
  \item The thesis title is italicized.
  \item The school name, address, and year follow, separated by commas.
\end{itemize}
Example output: Tullius T.~D., \emph{Title of PhD thesis}, University of Newcastle, UK, 1991.
